
\chapter{THẢO LUẬN VÀ SO SÁNH}

% ============================================
% 5.1. THẢO LUẬN TỔNG QUAN
% ============================================

\section{Thảo luận tổng quan kết quả thực nghiệm}

\subsection{Đánh giá chung về hiệu quả phân loại}

Trong quá trình thực hiện đề tài, nhóm đã triển khai
và đánh giá hai thuật toán phân loại là Multinomial
Naive Bayes và K-Nearest Neighbors (K-NN), kết hợp
với các phương pháp biểu diễn văn bản Bag-of-Words
(BoW) và TF-IDF.

Các thí nghiệm được thực hiện trên bộ dữ liệu gồm
1.223 phản hồi sinh viên, được chia thành 978 mẫu
huấn luyện và 245 mẫu kiểm thử.

Nhóm xây dựng một mô hình baseline sử dụng
BoW Unigram kết hợp Multinomial Naive Bayes với
tham số $\alpha=1,0$. Đồng thời, tám tổ hợp
biểu diễn--thuật toán được thử nghiệm và lựa chọn
siêu tham số bằng Stratified 3-Fold Cross-Validation.

Kết quả trình bày trong
Bảng~\ref{tab:ch4-ketqua-test}
cho thấy hiệu quả phân loại giữa các phương án
có sự khác biệt.

Mô hình TF-IDF Unigram kết hợp K-NN với $k=5$
được lựa chọn dựa trên kết quả kiểm định chéo,
đạt Macro F1-score trung bình cao nhất là 0,3647.

Tuy nhiên, khi đánh giá trên tập kiểm thử,
mô hình chỉ đạt Accuracy 0,3388 và Macro
F1-score 0,3369.

Trong khi đó, mô hình baseline BoW Unigram
kết hợp Multinomial Naive Bayes đạt Accuracy
0,3878 và Macro F1-score 0,3748.

Như vậy, Macro F1-score của mô hình được lựa chọn
thấp hơn baseline 0,0379 trên tập kiểm thử.

Kết quả này cho thấy mô hình có điểm kiểm định chéo
cao nhất chưa chắc đạt hiệu quả tốt hơn các phương
án khác trên một tập dữ liệu kiểm thử cụ thể.

Đây là một phát hiện quan trọng của thực nghiệm.
Nhóm không thay đổi mô hình được lựa chọn sau khi
quan sát kết quả test, nhằm giữ đúng nguyên tắc
đánh giá đã thiết lập trong quy trình học máy.


\subsection{Ý nghĩa của các chỉ số đánh giá}

Bên cạnh hai thuật toán học máy, nhóm sử dụng
DummyClassifier với chiến lược dự đoán nhãn
phổ biến nhất làm mốc tham khảo đơn giản.

Kết quả cho thấy DummyClassifier đạt Accuracy
0,4000 nhưng Macro F1-score chỉ đạt 0,1905.

Ngược lại, baseline Naive Bayes có Accuracy
thấp hơn, bằng 0,3878, nhưng Macro F1-score
cao hơn đáng kể, đạt 0,3748.

Nguyên nhân trực tiếp là DummyClassifier
luôn dự đoán cùng một lớp, dẫn đến khả năng
nhận diện các lớp còn lại rất hạn chế.

Trong trường hợp này, Accuracy chủ yếu phản
ánh tỷ lệ mẫu thuộc lớp được dự đoán thường xuyên,
chứ không cho biết mô hình có phân biệt tốt
cả ba nhóm cảm xúc hay không.

Do đó, việc chỉ sử dụng Accuracy có thể tạo
ra cách nhìn chưa đầy đủ về hiệu quả phân loại.

Với bài toán gồm ba lớp negative, neutral và
positive, nhóm sử dụng Macro F1-score làm tiêu
chí lựa chọn mô hình vì chỉ số này xem xét
đóng góp của từng lớp với trọng số như nhau.

Đồng thời, Precision, Recall và Confusion Matrix
được sử dụng để đánh giá chi tiết khả năng
nhận diện từng nhóm cảm xúc.

Như vậy, kết quả thực nghiệm cho thấy cần sử dụng
kết hợp nhiều chỉ số đánh giá, thay vì kết luận
chất lượng mô hình chỉ dựa trên một giá trị
độ chính xác chung.


% ============================================
% 5.2. SO SÁNH PHƯƠNG PHÁP VÀ THUẬT TOÁN
% ============================================

\section{So sánh các phương pháp biểu diễn và thuật toán}

\subsection{Ảnh hưởng của Bag-of-Words và TF-IDF}

Hai phương pháp biểu diễn văn bản được sử dụng
trong đề tài có cách xác định giá trị đặc trưng
khác nhau.

Bag-of-Words biểu diễn văn bản thông qua số lần
xuất hiện của từ hoặc cụm từ, trong khi TF-IDF
điều chỉnh trọng số đặc trưng dựa trên tần suất
xuất hiện trong văn bản và mức độ phổ biến của
từ trong tập huấn luyện.

Để phân tích ảnh hưởng của phương pháp biểu diễn,
nhóm so sánh Macro F1-score trên tập test khi
giữ nguyên thuật toán và cấu hình n-gram.

Kết quả được trình bày trong
Bảng~\ref{tab:ch5-bow-tfidf}.

\begin{table}[H]
\centering
\caption{So sánh Macro F1-score giữa BoW và TF-IDF}
\label{tab:ch5-bow-tfidf}
\renewcommand{\arraystretch}{1.15}
\small

\begin{tabularx}{\textwidth}{|X|r|r|r|}
\hline
\textbf{Cấu hình}
& \textbf{BoW}
& \textbf{TF-IDF}
& \textbf{Chênh lệch} \\
\hline

Unigram + NB
& 0,3702 & 0,3614 & -0,0088 \\
\hline

Unigram + KNN
& 0,3588 & 0,3369 & -0,0219 \\
\hline

Unigram--Bigram + NB
& 0,3436 & 0,3316 & -0,0120 \\
\hline

Unigram--Bigram + KNN
& 0,3271 & 0,3316 & +0,0045 \\
\hline

\end{tabularx}
\end{table}

\noindent
\textit{Ghi chú:} Chênh lệch được tính bằng
Macro F1 của TF-IDF trừ Macro F1 của BoW
trên cùng cấu hình thuật toán và n-gram.

Từ Bảng~\ref{tab:ch5-bow-tfidf}, có thể nhận thấy
BoW cho kết quả cao hơn TF-IDF ở ba trong bốn
trường hợp so sánh.

Đối với Naive Bayes kết hợp unigram, BoW đạt
Macro F1-score 0,3702, cao hơn TF-IDF với
kết quả 0,3614.

Tương tự, khi sử dụng K-NN với unigram,
BoW đạt Macro F1-score 0,3588, trong khi
TF-IDF đạt 0,3369.

Riêng cấu hình unigram kết hợp bigram với K-NN,
TF-IDF đạt Macro F1-score 0,3316, cao hơn
BoW với kết quả 0,3271. Tuy nhiên, mức
chênh lệch chỉ là 0,0045.

Những kết quả trên cho thấy TF-IDF không
mặc nhiên mang lại hiệu quả phân loại cao
hơn BoW trên bộ dữ liệu hiện tại.

Một giả thuyết có thể xem xét là việc điều chỉnh
trọng số từ bằng TF-IDF chưa tạo ra sự phân biệt
đủ rõ giữa các lớp cảm xúc.

Tuy nhiên, đây mới là giả thuyết dựa trên kết quả
thực nghiệm. Nhóm chưa thực hiện phân tích sâu
đóng góp của từng đặc trưng để khẳng định nguyên
nhân trực tiếp.

Vì vậy, kết luận chỉ giới hạn trong bộ dữ liệu,
các thuật toán và cấu hình đã thử nghiệm.


\subsection{Ảnh hưởng của unigram và bigram}

Ngoài phương pháp biểu diễn, nhóm tiến hành
so sánh các cấu hình unigram và unigram kết
hợp bigram.

Unigram biểu diễn văn bản thông qua từng từ
riêng lẻ, còn bigram bổ sung thông tin về
những cặp từ liên tiếp.

Việc sử dụng bigram có thể giúp mô hình nhận
diện một số cụm từ mang thông tin cảm xúc,
chẳng hạn như ``not good''.

Tuy nhiên, bổ sung bigram cũng làm tăng
số lượng đặc trưng tiềm năng và có thể khiến
ma trận biểu diễn văn bản trở nên thưa hơn.

Kết quả thực nghiệm cho thấy:

\begin{itemize}

    \item Với BoW kết hợp Naive Bayes,
    Macro F1-score giảm từ 0,3702 xuống
    0,3436 khi bổ sung bigram.

    \item Với TF-IDF kết hợp Naive Bayes,
    Macro F1-score giảm từ 0,3614 xuống
    0,3316.

    \item Với BoW kết hợp K-NN,
    Macro F1-score giảm từ 0,3588 xuống
    0,3271.

    \item Với TF-IDF kết hợp K-NN,
    Macro F1-score giảm từ 0,3369 xuống
    0,3316.

\end{itemize}

Như vậy, trong cả bốn trường hợp trên,
việc bổ sung bigram không cải thiện
Macro F1-score trên tập test hiện tại.

Một nguyên nhân có thể xảy ra là các cụm
hai từ xuất hiện không đủ thường xuyên
để cung cấp những tín hiệu phân loại ổn định.

Bên cạnh đó, một số phản hồi trong bộ dữ liệu
có cấu trúc ngữ nghĩa chưa rõ ràng. Vì vậy,
những cặp từ được bổ sung có thể chưa thể hiện
được các đặc trưng cảm xúc có ích cho mô hình.

Tuy nhiên, nhóm chưa thực hiện thí nghiệm
loại bỏ từng nhóm đặc trưng để xác nhận
các giả thuyết trên.

Kết quả cho thấy việc tăng số lượng đặc trưng
không đồng nghĩa với việc tăng hiệu quả phân loại.


\subsection{So sánh Naive Bayes và K-NN}

Hai thuật toán được lựa chọn trong đề tài
có nguyên lý hoạt động khác nhau.

Multinomial Naive Bayes là thuật toán phân loại
dựa trên xác suất, sử dụng giả định độc lập có
điều kiện giữa các đặc trưng khi đã biết lớp.

Trong khi đó, K-NN dự đoán nhãn dựa trên các
mẫu huấn luyện gần nhất theo một độ đo khoảng cách.

Đối với bài toán phân loại văn bản, Naive Bayes
có ưu điểm là cách huấn luyện đơn giản, tốc độ
xử lý tương đối nhanh và phù hợp với dữ liệu
được biểu diễn bằng tần suất hoặc trọng số từ.

Tuy nhiên, giả định độc lập có điều kiện
có thể chưa phản ánh đầy đủ mối quan hệ
giữa các từ trong câu.

K-NN có ưu điểm là nguyên lý trực quan,
dễ giải thích và không yêu cầu xây dựng
một hàm phân loại tham số phức tạp.

Tuy nhiên, khi dữ liệu văn bản được biểu diễn
trong không gian đặc trưng có số chiều lớn,
việc xác định những mẫu gần nhau về mặt vector
không nhất thiết đồng nghĩa với việc các
văn bản có cảm xúc giống nhau.

Trong kết quả thực nghiệm, Naive Bayes đạt
Macro F1-score cao hơn K-NN ở ba trong bốn
cấu hình biểu diễn tương ứng, còn một
cấu hình có kết quả bằng nhau khi làm tròn
đến bốn chữ số thập phân.

Đặc biệt, baseline BoW Unigram kết hợp
Naive Bayes đạt Macro F1-score 0,3748,
cao hơn mô hình TF-IDF Unigram kết hợp
K-NN được lựa chọn bằng Cross-Validation.

Điều này cho thấy Naive Bayes là một phương án
đáng chú ý trong phạm vi bộ dữ liệu và
các cấu hình đã thử nghiệm.

Tuy nhiên, nhóm không kết luận Naive Bayes
luôn tốt hơn K-NN đối với mọi bài toán
phân loại văn bản.

Hiệu quả thực tế còn phụ thuộc vào dữ liệu,
phương pháp biểu diễn, siêu tham số và
cách thiết kế thí nghiệm.


% ============================================
% 5.3. KHÁC BIỆT GIỮA CV VÀ TEST
% ============================================

\section{Phân tích sự khác biệt giữa Cross-Validation và tập kiểm thử}

\subsection{Chênh lệch kết quả CV và test}

Trong giai đoạn lựa chọn mô hình, nhóm sử dụng
Stratified 3-Fold Cross-Validation trên tập
huấn luyện với Macro F1-score làm tiêu chí.

Kết quả cho thấy TF-IDF Unigram kết hợp K-NN
với $k=5$ đạt Macro F1 trung bình cao nhất,
bằng 0,3647.

Cấu hình BoW Unigram kết hợp Naive Bayes
với $\alpha=0,5$ đạt điểm CV là 0,3637,
chỉ thấp hơn khoảng 0,0010.

Mức chênh lệch này tương đối nhỏ.

Ngoài ra, độ lệch chuẩn điểm CV của hai cấu
hình lần lượt xấp xỉ 0,0088 và 0,0148.

Điều này cho thấy mức chênh lệch điểm trung
bình giữa hai cấu hình nhỏ hơn mức biến thiên
quan sát được giữa các fold.

Tuy nhiên, chỉ dựa vào ba fold và các độ lệch
chuẩn này chưa đủ để thực hiện kết luận
thống kê về sự khác biệt giữa hai mô hình.

Khi đánh giá trên tập test, mô hình TF-IDF
Unigram kết hợp K-NN đạt Macro F1-score
0,3369, thấp hơn điểm CV trung bình khoảng
0,0278.

Trong khi đó, baseline BoW Unigram kết hợp
Naive Bayes với $\alpha=1,0$ đạt Macro
F1-score 0,3748 trên test.

Các kết quả cho thấy cấu hình được lựa chọn
bằng CV không thể hiện được ưu thế trên
tập kiểm thử hiện tại.

Cần lưu ý rằng baseline sử dụng $\alpha=1,0$,
trong khi cấu hình BoW Unigram kết hợp Naive Bayes
được chọn qua CV sử dụng $\alpha=0,5$.

Vì vậy, khi so sánh hai phương án này, nhóm
phải phân biệt rõ giữa baseline cố định và
mô hình có siêu tham số được tinh chỉnh.


\subsection{Những nguyên nhân có thể ảnh hưởng đến kết quả}

Sự khác biệt giữa điểm CV và điểm test có thể
liên quan đến nhiều yếu tố.

Thứ nhất, số lượng mẫu huấn luyện và kiểm thử
còn giới hạn. Kết quả trên một lần chia
Train/Test có thể chịu ảnh hưởng từ những
mẫu cụ thể được đưa vào mỗi tập.

Thứ hai, nội dung văn bản trong bộ dữ liệu
có thể chưa cung cấp đủ thông tin để phân
biệt rõ ba nhãn cảm xúc.

Thứ ba, các phương pháp BoW và TF-IDF chủ yếu
dựa trên thông tin từ hoặc cụm từ, trong khi
cảm xúc có thể phụ thuộc vào ngữ cảnh
và ý nghĩa tổng thể của câu.

Thứ tư, những giá trị siêu tham số được thử
nghiệm mới nằm trong một phạm vi tương đối
nhỏ, gồm ba giá trị $\alpha$ đối với Naive Bayes
và ba giá trị $k$ đối với K-NN.

Các yếu tố trên chỉ là những giả thuyết
cần được kiểm chứng thông qua các thí nghiệm
bổ sung.

Không thể khẳng định mô hình bị overfitting
chỉ dựa trên việc điểm test thấp hơn
điểm CV ở một lần đánh giá.

Để đánh giá độ ổn định tốt hơn, nhóm có thể
thực hiện thêm các lần chia dữ liệu khác
hoặc kiểm định chéo lặp lại trong các
nghiên cứu tiếp theo.

Trong phạm vi bài tập lớn hiện tại, nhóm
giữ nguyên mô hình đã được lựa chọn bằng CV
và không sử dụng kết quả test để điều chỉnh
lại lựa chọn sau khi xem điểm số.


% ============================================
% 5.4. PHÂN TÍCH DỰ ĐOÁN SAI
% ============================================

\section{Phân tích các trường hợp dự đoán sai}

\subsection{Phân tích các lớp thường bị nhầm lẫn}

Kết quả đánh giá mô hình TF-IDF Unigram
kết hợp K-NN cho thấy có 162 trường hợp
dự đoán sai trong tổng số 245 phản hồi
kiểm thử.

Ma trận nhầm lẫn được trình bày tại
Bảng~\ref{tab:ch4-confusion-matrix}
và Hình~\ref{fig:ch4-cm-knn}.

Trong đó, cặp nhãn thường bị nhầm lẫn
nhiều nhất là positive được dự đoán thành
negative, với 41 trường hợp.

Ngoài ra, có 29 trường hợp neutral bị dự
đoán thành negative và 27 trường hợp
positive bị dự đoán thành neutral.

Đối với nhãn negative, mô hình dự đoán
đúng 31 trên 73 mẫu, tương ứng Recall
0,4247.

Đối với nhãn neutral, mô hình dự đoán
đúng 22 trên 74 mẫu, đạt Recall 0,2973.

Đối với nhãn positive, mô hình dự đoán
đúng 30 trên 98 mẫu, đạt Recall 0,3061.

Như vậy, cả ba lớp đều có Recall dưới 0,5.
Điều này cho thấy mô hình chưa nhận diện
chính xác được phần lớn các phản hồi
thuộc từng lớp cảm xúc.

Đặc biệt, lớp neutral có Recall thấp nhất.
Đây là lớp cần được chú ý khi phân tích lỗi,
bởi mô hình thường phân loại một số phản
hồi trung tính sang positive hoặc negative.

Kết quả cũng cho thấy việc phân biệt
hai nhóm cảm xúc đối lập positive và negative
chưa đạt hiệu quả tốt trên bộ dữ liệu này.


\subsection{Phân tích một số phản hồi cụ thể}

Để tìm hiểu đặc điểm của các trường hợp
dự đoán sai, nhóm xem xét các phản hồi
được xuất từ notebook
\texttt{04\_Evaluation.ipynb}.

Các phản hồi dưới đây là những ví dụ
có trong tập kiểm thử thực tế.

\textbf{Trường hợp 1: Phản hồi positive bị dự
đoán thành negative.}

\begin{quote}
\textit{``Bad everything rise through similar
kid talk item day.''}
\end{quote}

Nhãn được cung cấp trong dữ liệu là
\texttt{positive}, trong khi mô hình
dự đoán \texttt{negative}.

Trong câu xuất hiện từ ``Bad'', thường
mang sắc thái tiêu cực khi đứng trong
một ngữ cảnh phù hợp.

Tuy nhiên, toàn bộ câu chưa hình thành
một nhận xét rõ ràng về trải nghiệm học tập.

Do đó, chưa thể kết luận mô hình dự đoán
negative chỉ vì sự xuất hiện của từ ``Bad''.

Trường hợp này cho thấy cần xem xét
tính rõ nghĩa của văn bản và sự phù hợp
giữa nội dung với nhãn được cung cấp.


\textbf{Trường hợp 2: Phản hồi neutral bị dự
đoán thành negative.}

\begin{quote}
\textit{``Trade leader option interview cultural
center one seek activity take push forget detail any.''}
\end{quote}

Nhãn được cung cấp trong dữ liệu là
\texttt{neutral}, còn mô hình dự đoán
\texttt{negative}.

Phản hồi chứa nhiều từ và cụm từ có liên hệ
ngữ nghĩa chưa rõ ràng.

Nội dung không thể hiện một đánh giá
trực tiếp về hoạt động học tập, môn học
hoặc trải nghiệm của sinh viên.

Vì vậy, việc xác định sắc thái cảm xúc
chỉ từ nội dung câu có thể gặp khó khăn.

Đây là một dấu hiệu cho thấy mô hình
có thể thiếu thông tin ngữ cảnh phù hợp
để phân biệt các lớp cảm xúc.


\textbf{Trường hợp 3: Phản hồi neutral bị dự
đoán thành positive.}

\begin{quote}
\textit{``Feeling some four room shoulder
challenge join development wind key poor black.''}
\end{quote}

Nhãn được cung cấp trong dữ liệu là
\texttt{neutral}, nhưng mô hình dự đoán
\texttt{positive}.

Trong câu có một số từ có thể liên quan
đến cảm xúc hoặc đánh giá như ``Feeling'',
``challenge'' và ``poor''.

Tuy nhiên, các từ này không hình thành
một nội dung phản hồi có ý nghĩa rõ ràng
khi xét trong toàn bộ câu.

Vì vậy, chưa thể xác định chính xác
lý do khiến mô hình đưa ra dự đoán
positive thay vì neutral.

Trường hợp này cần được xem xét cùng
những ví dụ khác trước khi đưa ra
kết luận về nguyên nhân lỗi.


\subsection{Nhận xét chung từ phân tích lỗi}

Qua ba trường hợp được lựa chọn, nhóm
nhận thấy một số phản hồi có đặc điểm
chung là cấu trúc câu thiếu tự nhiên,
mối liên hệ giữa các từ chưa rõ ràng
và không biểu đạt trực tiếp một thái độ
đối với trải nghiệm học tập.

Đặc điểm này có thể gây khó khăn cho
việc xác định cảm xúc dựa trên nội dung
văn bản, kể cả khi sử dụng các phương
pháp biểu diễn đặc trưng khác nhau.

Tuy nhiên, ba ví dụ trên không đủ
đại diện cho toàn bộ các trường hợp
dự đoán sai.

Bên cạnh đó, nhãn do bộ dữ liệu cung cấp
được sử dụng làm nhãn tham chiếu trong
đánh giá mô hình, nhưng nhóm chưa thực
hiện quy trình gán nhãn lại độc lập.

Vì vậy, nhóm không khẳng định những
nhãn trên là sai, cũng không kết luận
chất lượng ngữ nghĩa là nguyên nhân duy
nhất khiến mô hình có kết quả thấp.

Để đánh giá đầy đủ hơn, cần xem xét
nhiều trường hợp dự đoán sai theo từng
cặp nhãn và đối chiếu với quy tắc
gán nhãn của nguồn dữ liệu.


% ============================================
% 5.5. HẠN CHẾ
% ============================================

\section{Hạn chế của dữ liệu và giải pháp hiện tại}

\subsection{Hạn chế về dữ liệu}

Bộ dữ liệu được sử dụng có 1.223 phản hồi,
đáp ứng phạm vi thực nghiệm cơ bản của
học phần Nhập môn Học máy.

Tuy nhiên, quy mô này vẫn còn giới hạn
nếu muốn đánh giá khả năng ứng dụng
trên nhiều nhóm phản hồi khác nhau.

Bên cạnh đó, một số phản hồi có nội dung
chưa rõ nghĩa và thiếu ngữ cảnh học tập
cụ thể.

Đây là vấn đề cần được tiếp tục xem xét
vì nhãn cảm xúc được dự đoán chỉ dựa
trên nội dung văn bản.

Ngoài ra, mặc dù dữ liệu không có
giá trị thiếu hoặc trùng lặp theo
các phép kiểm tra đã thực hiện,
nhóm chưa có đầy đủ bằng chứng để
xác nhận tính chính xác của tất cả
các nhãn được cung cấp.

Do đó, việc kiểm tra cấu trúc dữ liệu
chưa đủ để bảo đảm chất lượng ngữ nghĩa
phục vụ bài toán phân loại cảm xúc.


\subsection{Hạn chế của phương pháp tiền xử lý}

Quy trình tiền xử lý hiện tại thực hiện
chuẩn hóa chữ thường, xử lý URL,
dấu câu, ký tự đặc biệt và khoảng
trắng thừa.

Nhóm giữ lại stopword, chữ số và các
từ phủ định nhằm hạn chế việc loại
bỏ thông tin cảm xúc một cách máy móc.

Tuy nhiên, quy tắc làm sạch vẫn có thể
làm mất một số tín hiệu được thể hiện
thông qua dấu câu hoặc cách trình bày
văn bản.

Bên cạnh đó, nhóm chưa thực hiện
các thí nghiệm có kiểm soát để so
sánh trực tiếp hiệu quả của nhiều
quy tắc tiền xử lý khác nhau.

Vì vậy, chưa thể kết luận quy trình
làm sạch hiện tại là phương án tối ưu
đối với bộ dữ liệu.


\subsection{Hạn chế của phương pháp biểu diễn}

Bag-of-Words và TF-IDF có ưu điểm
là dễ triển khai, dễ giải thích và
phù hợp với các thuật toán học máy
cơ bản.

Tuy nhiên, hai phương pháp này
chủ yếu biểu diễn văn bản dựa
trên sự xuất hiện hoặc trọng số
của từ và cụm từ.

Do đó, chúng chưa thể hiện đầy đủ
ý nghĩa tổng thể, thứ tự từ trong
các cấu trúc dài hoặc quan hệ
ngữ nghĩa phức tạp của câu.

Việc bổ sung bigram có thể giữ
được một phần thông tin từ hai
từ liên tiếp nhưng chưa giải quyết
hoàn toàn những hạn chế trên.

Đối với các phản hồi thiếu ngữ cảnh,
những hạn chế về biểu diễn có thể
khiến mô hình khó xác định chính
xác nhãn cảm xúc.


\subsection{Hạn chế trong quá trình thực nghiệm}

Nhóm đã sử dụng Cross-Validation
để lựa chọn siêu tham số trên tập
huấn luyện và giữ riêng tập test
cho bước đánh giá cuối cùng.

Tuy nhiên, quá trình đánh giá
hiện tại mới sử dụng một lần
chia Train/Test cố định.

Do đó, kết quả chưa phản ánh
đầy đủ sự biến thiên có thể xảy
ra khi thay đổi cách chia dữ liệu.

Ngoài ra, phạm vi thử nghiệm
siêu tham số còn tương đối hẹp.

Nhóm mới xem xét ba giá trị $\alpha$
đối với Naive Bayes và ba giá trị
$k$ đối với K-NN.

Vì vậy, chưa có cơ sở để khẳng
định những giá trị được chọn
là tối ưu trên mọi cấu hình
hoặc trên các bộ dữ liệu khác.

Cuối cùng, những kết quả đánh giá
hiện tại còn thấp và chưa chứng
minh được hiệu quả đủ tin cậy
để triển khai hệ thống trong
môi trường ứng dụng thực tế.


% ============================================
% 5.6. ĐỀ XUẤT CẢI TIẾN
% ============================================

\section{Đề xuất hướng cải tiến và khả năng ứng dụng}

\subsection{Cải thiện chất lượng dữ liệu}

Một hướng cải tiến quan trọng là
kiểm tra và nâng cao chất lượng
của bộ dữ liệu đầu vào.

Nhóm có thể xây dựng quy tắc
đánh giá nhãn rõ ràng, sau đó
kiểm tra thủ công một tập mẫu
đại diện cho ba lớp cảm xúc.

Đặc biệt, nên ưu tiên kiểm tra
những phản hồi có nội dung
khó hiểu hoặc có sự không
thống nhất rõ ràng giữa
văn bản và nhãn được cung cấp.

Nếu có điều kiện, nhóm có thể
bổ sung những phản hồi thực tế
có nội dung rõ nghĩa và được
gán nhãn theo quy trình thống nhất.

Tuy nhiên, việc thay đổi hoặc
thay thế bộ dữ liệu trong phạm
vi bài tập lớn cần được sự
đồng ý của giảng viên theo
yêu cầu của học phần.


\subsection{Mở rộng thử nghiệm tiền xử lý và đặc trưng}

Trong các nghiên cứu tiếp theo,
nhóm có thể thử nghiệm thêm
những phương án xử lý văn bản
và so sánh kết quả theo cùng
một quy trình đánh giá.

Chẳng hạn, nhóm có thể phân tích
tác động của việc giữ hoặc loại
bỏ một số dấu câu, điều chỉnh
ngưỡng \texttt{min\_df},
\texttt{max\_df} hoặc số lượng
đặc trưng tối đa.

Các thí nghiệm cần được thiết
kế sao cho chỉ thay đổi một
hoặc một nhóm tham số có
lý do, nhằm thuận tiện cho
việc phân tích tác động.

Bên cạnh đó, nhóm có thể
tiếp tục khảo sát những
cấu hình n-gram phù hợp
với đặc điểm văn bản.

Việc mở rộng phương pháp
biểu diễn cần được thực
hiện có kiểm soát, thay vì
mặc định số lượng đặc
trưng càng lớn thì kết
quả càng tốt.


\subsection{Đánh giá độ ổn định của mô hình}

Để đánh giá khả năng khái
quát hóa tốt hơn, nhóm
có thể thực hiện thêm các
thí nghiệm với nhiều
giá trị random_state
hoặc sử dụng kiểm định
chéo lặp lại trên dữ liệu
huấn luyện.

Việc sử dụng nhiều lần
đánh giá giúp quan sát
mức độ biến thiên của
các chỉ số khi dữ liệu
được phân chia khác nhau.

Nhóm cũng có thể mở rộng
phạm vi thử nghiệm tham
số $\alpha$ và $k$,
nhưng việc lựa chọn
tham số vẫn cần được
thực hiện bằng dữ liệu
huấn luyện hoặc validation.

Tập test dùng để đánh
giá cuối cùng không
nên được sử dụng lặp
đi lặp lại để lựa
chọn mô hình.

Nếu có thêm dữ liệu,
một tập kiểm thử độc
lập mới sẽ giúp đánh
giá hiệu quả thực tế
của giải pháp một
cách đáng tin cậy hơn.


\subsection{Khả năng ứng dụng thực tế}

Mục tiêu dài hạn của
bài toán phân loại cảm
xúc là hỗ trợ quá trình
tổng hợp phản hồi
của sinh viên.

Một hệ thống có chất
lượng dự đoán tốt có
thể hỗ trợ thống kê
số lượng phản hồi
tích cực, trung tính
và tiêu cực.

Từ đó, nhà trường
hoặc giảng viên có
thể sử dụng kết quả
như một nguồn thông
tin tham khảo trong
quá trình xem xét
ý kiến người học.

Tuy nhiên, các kết
quả thực nghiệm hiện
tại chưa chứng minh
mô hình đủ chính xác
để ứng dụng trực tiếp.

Đặc biệt, những
trường hợp nhầm
positive thành
negative có thể
dẫn đến việc diễn
giải không đúng
sắc thái phản hồi
nếu kết quả tự
động được sử dụng
mà không kiểm tra.

Vì vậy, trong
giai đoạn hiện tại,
mô hình chỉ nên
được xem là sản
phẩm thực nghiệm
phục vụ mục đích
học tập và nghiên
cứu.

Nếu tiếp tục phát
triển thành ứng
dụng thực tế,
cần cải thiện
chất lượng dữ liệu,
đánh giá lại mô hình
và có cơ chế để
con người kiểm tra
những phản hồi
khó phân loại.


% ============================================
% 5.7. TỔNG KẾT CHƯƠNG
% ============================================

\section{Tổng kết chương}

Chương 5 đã thảo luận
và so sánh hiệu quả
của các phương pháp
biểu diễn văn bản
và thuật toán học máy
được sử dụng trong
đề tài.

Kết quả thực nghiệm
cho thấy Bag-of-Words
không nhất thiết kém
hơn TF-IDF, và việc
bổ sung bigram chưa
mang lại sự cải thiện
Macro F1-score trên
tập kiểm thử hiện tại.

Bên cạnh đó, mô hình
TF-IDF Unigram kết hợp
K-NN được lựa chọn bằng
Cross-Validation có
Macro F1-score thấp hơn
baseline BoW Unigram
kết hợp Naive Bayes
trên tập test.

Qua phân tích ma trận
nhầm lẫn và một số
trường hợp dự đoán sai,
nhóm nhận thấy mô hình
còn gặp khó khăn khi
phân biệt các lớp
cảm xúc.

Một số phản hồi có
nội dung thiếu rõ
nghĩa, đồng thời
các phương pháp
biểu diễn hiện tại
chưa thể hiện đầy
đủ thông tin ngữ
nghĩa của văn bản.

Tuy nhiên, những
hạn chế này mới
được xem xét dưới
dạng giả thuyết
dựa trên kết quả
thực nghiệm và
các mẫu quan sát.

Nhóm chưa có đủ
bằng chứng để khẳng
định một nguyên nhân
duy nhất dẫn đến
hiệu quả phân loại
thấp.

Những kết quả và
nhận xét trong
chương là cơ sở
để tổng kết quá
trình thực hiện,
đánh giá mức độ
đạt được mục tiêu
và đề xuất hướng
phát triển tiếp
theo trong Chương 6.
